\section{Métricas de conversión en producto y ROI}

\subsection{Lista de métricas}

Después de poner en producción el sistema DSPy es necesario recopilar de manera simultánea las siguientes métricas: costo de LLM por consulta, tasa de éxito de invocación de herramientas, tiempo de ejecución del optimizador, cobertura de automatización (proporción de agent) y la calificación de satisfacción del usuario.

\begin{table}[H]
\centering
\caption{Métricas para medir el ROI}
\begin{tabular}{p{0.22\textwidth}p{0.28\textwidth}p{0.28\textwidth}}
\toprule
Métrica & Ventana de observación & Dirección esperada \\
\midrule
Costo / consulta & Diaria & Reducir 10--20\% \\
Cobertura de automatización del agent & Semanal & Elevar a >70\% de la carga de trabajo \\
Satisfacción del usuario & Cada release & Elevar 0.2 puntos (sobre 5) \\
Éxito de invocación de herramientas & En tiempo real & > 95\% \\
\bottomrule
\end{tabular}
\end{table}

\subsection{Explicación del ROI}

El ROI no solo considera el ahorro en costos de la API, sino que también debe tomar en cuenta la reducción del tiempo del desarrollador. Mediante el análisis de trace de la Evidence Matrix es posible cuantificar el aumento en la velocidad de respuesta que `optimizer improvement` aporta al negocio y, a partir de ahí, calcular `value delivered / cost spent`.

\subsection{Resumen de la sección}

Las métricas de conversión en producto aseguran que la inversión en DSPy esté alineada con los objetivos del negocio; el cálculo del ROI requiere traducir las mejoras técnicas en cobertura del agent y en satisfacción del usuario.
